\documentclass[11pt, letterpaper]{article}
\usepackage[utf8]{inputenc}
\usepackage[T1]{fontenc}
\usepackage{amsmath}
\usepackage{amsfonts}
\usepackage{amssymb}
\usepackage{geometry}
\usepackage{verbatim}

\geometry{letterpaper, margin=1in}

\title{\textbf{Problem B: Sayaka's Blend}}
\author{Time Limit: 2000ms | Memory Limit: 256MB}
\date{}

\begin{document}

\maketitle

\section*{Description}

"I prefer my coffee precise. Just like my life."

Sayaka Saeki is experimenting with coffee blends at the cafe. She has $N$ distinct types of coffee beans available.
The $i$-th bean has a \textbf{Flavor Score} $F_i$ and a \textbf{Caffeine Content} $C_i$.

Sayaka wants to create a blend by selecting a \textbf{non-empty} subset of beans. The blend must meet a strict requirement: the \textbf{Average Caffeine Content} of the selected beans must be exactly $K$.
Among all valid blends, she wants to maximize the \textbf{Total Flavor Score} ($\sum F_i$).

\section*{Input}

The first line contains two integers $N$ and $K$ ($1 \le N \le 40$, $1 \le K \le 10^9$).
The next $N$ lines each contain two integers $F_i$ and $C_i$ ($1 \le F_i, C_i \le 10^9$).

\section*{Output}

Output the maximum total flavor score possible. If no subset has an average caffeine content of exactly $K$, output 0.

\section*{Examples}

\subsection*{Example 1}
\textbf{Input:}
\begin{verbatim}
4 10
100 10
50 20
60 5
40 5
\end{verbatim}

\textbf{Output:}
\begin{verbatim}
250
\end{verbatim}

\textbf{Explanation:}
We need the average caffeine to be exactly 10.
If we select Bean 1 ($C=10$), Bean 2 ($C=20$), Bean 3 ($C=5$), and Bean 4 ($C=5$):
Sum of Caffeine $= 10 + 20 + 5 + 5 = 40$.
Count $= 4$.
Average $= 40 / 4 = 10$. This is valid.
Total Flavor $= 100 + 50 + 60 + 40 = 250$.

\subsection*{Example 2}
\textbf{Input:}
\begin{verbatim}
3 100
10 10
20 20
30 30
\end{verbatim}

\textbf{Output:}
\begin{verbatim}
0
\end{verbatim}

\end{document}
